\scanpage{169}
\uline{Application} Un schéma étale \inserted{$T$ \uline{sur $S$}} est dit à \uline{groupe}
\uline{structural} $\pi$, \uline{fibre $F$}, si on a
\deleted{\unsure{si on se donne un} sous-schéma de $\underline{\operatorname{Aut}}_S(T)$, qui est}
donné \unsure{une opération fidèle|des opérations fidèles} de $\pi$ sur $F$,
et un sous-fibré \inserted{$\pi_S$} principal de
$\underline{\operatorname{Isom}}_S(E_S,T)$. Si $S$ est connexe et
muni du pt géométrique $\xi$, il revient
au $=$ de se donner un $T$ avec
une telle structure, ou de se
donner des opérations \deleted{continues}
de $\widetilde{\pi}_1(S,\xi)$ \deleted{\unsure{à travers}} sur $F$ à
travers $\widetilde{\pi}_1(S,\xi)\unsure{:|\,}\to\deleted{\unsure{G}}\pi$.

\uline{Exemple} Schémas en groupes comm.
\unsure{constants} tordus sur $S$, \deleted{à fibres \unsure{gé.}} \inserted{de}
génération finie sur $S$, correspondent
\deleted{\unsure{aux|au}} $(\widetilde{\pi}_1(S,\xi))$-Modules admissibles
\deleted{de \unsure{présentat|présentation}} qui sont des $\mathbb Z$-modules
de type fini \unsure{---}
% archive sheet number: 168
